\documentclass{article}

%**************************************************************
% Importazione package
%**************************************************************

% permette di modificare i margini
%\usepackage[top=3.1cm, bottom=3.1cm, left=2.2cm, right=2.2cm]{geometry}
\usepackage[a4paper]{geometry}

% necessario per risolvere il problema del carattere invisibile per l'a capo
% \DeclareUnicodeCharacter{00A0}{ }

% per scrivere in italiano e in inglese;
% l'ultima lingua (l'italiano) risulta predefinita
\usepackage[english, italian]{babel}

% imposta lo stile italiano per i paragrafi
\usepackage{parskip}

% fornice elenchi numerati in ordine inverso
\usepackage{etaremune}

% comandi per l'appendice
\usepackage{appendix}
\renewcommand\appendixtocname{Appendici}
\renewcommand{\appendixpagename}{Appendici}

% numera anche i paragrafi
\setcounter{secnumdepth}{6}

% elenca anche i paragrafi nell'indice
\setcounter{tocdepth}{6}

% permetti di definire dei colori
\usepackage[usenames,dvipsnames]{color}

% permette di inserire le immagini/tabelle esattamente dove viene usato il
% comando \begin{figure}[H] ... \end{figure}
% evitando che venga spostato in automatico
\usepackage{float}

% permette l'inserimento di url e di altri tipi di collegamento
\usepackage[colorlinks=true]{hyperref}

\hypersetup{
    colorlinks=true, % false: boxed links; true: colored links
    citecolor=black,
    filecolor=black,
    linkcolor=black, % color of internal links
    urlcolor=Maroon  % color of external links
}

% immagini
\usepackage{graphicx}

% permette di riferirsi all'ultima pagina con \pageref{LastPage}
\usepackage{lastpage}

% tabelle su più pagine
\usepackage{longtable}

% per avere dei comandi in più da poter usare sulle tabelle
\usepackage{booktabs}

% tabelle con il campo X per riempire lo spazio rimanente sulla riga
\usepackage{tabularx}

% multirow per tabelle
\usepackage{multirow}

% colore di sfondo per le celle
\usepackage[table]{xcolor}
%\rowcolors{2}{gray!25}{} % Colora righe alterne in grigio (non funziona bene)

% permette di fare longtable larghe tutta la pagina (parametro x)
% su Ubuntu non si può installare il pacchetto, deve essere in model/
\usepackage{tabu}

% imposta lo spazio tra le righe di una tabella
\setlength{\tabulinesep}{6pt}

% definisci un nuovo tipo di colonna P che permette di andare a capo con \newline
% e giustificata a sinistra
\usepackage{array}
\usepackage{ragged2e}
\newcolumntype{P}[1]{>{\RaggedRight \hspace{0pt}}m{#1}}

% personalizza l'intestazione e piè di pagina
\usepackage{fancyhdr}

% permette di inserire caratteri speciali
\usepackage{textcomp}

% permette di aggiustare i margini e centrare tabelle e figure
\usepackage{changepage}

%Permette di includere i grafici a barre
%IMPORTANTE: deve essere caricato prima di /pgfgantt altrimenti causa conflitto
\usepackage{pgfplots}

% permette di includere i diagrammi Gantt
% su Ubuntu non si può installare il pacchetto, deve essere in model/
\usepackage{pgfgantt}

% permette di includere i grafici a torta
\usepackage{pgf-pie}

% necessario per pgf-pie
\usepackage{tikz}

% permette i path delle immagini con gli spazi
\usepackage{grffile}

% ruota le immagini
\usepackage{rotating}

% permetti di calcolare le larghezze facendo calcoli
\usepackage{calc}


%%
%% Definizione di path globali
%%
% Path assoluta a questa directory, va passato come argomento il resto della path
\newcommand{\ModelPath}[1]{/data/sources/model/#1}

% Path per accedere a modelassets
\newcommand{\ModelAssets}[1]{\ModelPath{modelassets}/#1}

% Path globali per la ricerca di immagini
\graphicspath{
	{/data/logo/PNG}
	{/data/logo/Square}
}



\fancypagestyle{plain}{
	% cancella tutti i campi di intestazione e piè di pagina
	\fancyhf{}
	\lhead{
		\includegraphics[width=1.5cm, keepaspectratio=true]{argo_icona.png}
		\parbox[b]{5cm}{
			\emph{\GroupName{}} \vspace{0pt} \\
			\emph{Progetto \ProjectName{}} \vspace{0pt}
		}
	}
	\chead{}


	\lfoot{
		\DocTitle{} \\
		% differenzia a seconda che \DocVersion{} stampi testo o no
		\setbox0=\hbox{\DocVersion{}\unskip}\ifdim\wd0=0pt
			% nulla
		\else
			v \DocVersion{}
		\fi
	}
	\rfoot{\thepage{} di \pageref{LastPage}}

	% Visualizza una linea orizzontale in cima e in fondo alla pagina
	\renewcommand{\headrulewidth}{0.3pt}
	\renewcommand{\footrulewidth}{0.3pt}
}
\setlength{\headheight}{30pt}
\pagestyle{plain}

% allarga l'header a tutta la pagina
%\fancyhfoffset[L]{\oddsidemargin + \hoffset + 1in}
%\fancyhfoffset[R]{\evensidemargin + \marginparwidth - \marginparsep}

% Per inserire del codice sorgente formattato
\usepackage{listings}
\definecolor{darkgray}{rgb}{.4,.4,.4}
\definecolor{purple}{rgb}{0.65, 0.12, 0.82}
\definecolor{mygreen}{rgb}{0,0.6,0}
\definecolor{mygray}{rgb}{0.5,0.5,0.5}
\definecolor{mymauve}{rgb}{0.58,0,0.82}

\lstset{
  extendedchars=true,          % lets you use non-ASCII characters
  inputencoding=utf8,   % converte i caratteri utf8 in latin1, richiede \usepackage{listingsutf8} anzichè listings
  basicstyle=\ttfamily,        % the size of the fonts that are used for the code
  breakatwhitespace=false,     % sets if automatic breaks should only happen at whitespace
  breaklines=true,             % sets automatic line breaking
  captionpos=b,                % sets the caption-position to top
  commentstyle=\color{mygreen},   % comment style
  frame=none,               % adds a frame around the code
  keepspaces=true,            % keeps spaces in text, useful for keeping indentation of code (possibly needs columns=flexible)
  keywordstyle=\color{blue}\bfseries,     % keyword style
  numbers=none,               % where to put the line-numbers; possible values are (none, left, right)
  numbersep=5pt,              % how far the line-numbers are from the code
  numberstyle=\color{mygray}, % the style that is used for the line-numbers
  rulecolor=\color{black},    % if not set, the frame-color may be changed on line-breaks within not-black text (e.g. comments (green here))
  showspaces=false,           % show spaces everywhere adding particular underscores; it overrides 'showstringspaces'
  showstringspaces=false,     % underline spaces within strings only
  showtabs=false,             % show tabs within strings adding particular underscores
  stepnumber=5,               % the step between two line-numbers. If it's 1, each line will be numbered
  stringstyle=\color{red},    % string literal style
  tabsize=4,                  % sets default tabsize
  firstnumber=1      % visualizza i numeri dalla prima linea
}

% Permetti di utilizzare il grassetto per i caratteri Typewriter (per es. il font di \code{...} e \file{...})
\usepackage[T1]{fontenc}
\usepackage{lmodern}

% Permette di resettare il contatore delle note a piè di pagina ad ogni pagina
\usepackage[perpage, bottom, stable]{footmisc}
\patchcmd{\footref}{\ref}{\ref*}{}{}

% Per poter usare il carattere ° (anche nelle formule matematiche)
\usepackage{gensymb}
\usepackage{newunicodechar}
\newunicodechar{°}{\degree}

% Per importare svg
\usepackage{svg}
%\setsvg{inkscape = inkscape -z -D}

% Per numerare le tabelle e le figure con la sezione in cui si trovano
\usepackage{amsmath}
\numberwithin{figure}{section}
\numberwithin{table}{section}

% Per aggiungere un nuovo lovello di profondita (subsubparagraph)

\makeatletter
\newcounter{subsubparagraph}[subparagraph]
\renewcommand\thesubsubparagraph{%
	\thesubparagraph.\@arabic\c@subsubparagraph}
\newcommand\subsubparagraph{%
	\@startsection{subsubparagraph}    % counter
	{6}                              % level
	{\parindent}                     % indent
	{3.25ex \@plus 1ex \@minus .2ex} % beforeskip
	{-1em}                           % afterskip
	{\normalfont\normalsize\bfseries}}
\newcommand\l@subsubparagraph{\@dottedtocline{6}{12em}{5em}}
\newcommand{\subsubparagraphmark}[1]{}
\def\endsubsubparagraph{\par
	\if@noskipsec
	\ifx\@svsechd\@undefined\else\leavevmode\fi
	\fi}
\makeatletter
\def\toclevel@subsubparagraph{6}
\makeatother
% Use \ul{arg} to undlerline text
\usepackage{soul}

% Per le pagine in orientamento orizzontale
\usepackage{pdflscape}
\usepackage{afterpage}

% Set vertical space in tables
\def\arraystretch{1.5}

% Modifica font
\usepackage{fontspec}

\setmainfont{Poppins}[
    Path=\ModelAssets{poppins/} ,
    Extension = .otf ,
    UprightFont=*-Regular ,
    BoldFont=*-Bold ,
    ItalicFont=*-Italic ,
    BoldItalicFont=*-BoldItalic
]
%%%%%%%%%%%%%%
%  COSTANTI  %
%%%%%%%%%%%%%%

% In questa prima parte vanno definite le 'costanti' utilizzate da due o più documenti.

% Serve a dare la giusta formattazione alle parole presenti nel glossario
% il nome del comando \glossary è già usato da LaTeX
% permette di distingure la prima occorrenza della parola nel documento
\usepackage{xparse}
\ExplSyntaxOn
\seq_new:N \l_hernan_seq
\NewDocumentCommand{\glossario}{m}{
    \seq_if_in:NnTF{\l_hernan_seq}{#1}{
      \textit{#1\ped{\ped{G}}}
    }{
      \seq_put_left:Nn{\l_hernan_seq}{#1}
      \textit{#1\ped{\ped{G}}}
    }
}
\ExplSyntaxOff

% Meglio non mettere gli \emph dentro le costanti, in certi casi creano problemi
\newcommand{\GroupName}{Argo}
\newcommand{\GroupEmail}{argo.unipd@gmail.com}
\newcommand{\ProjectName}{ChatSQL}
\newcommand{\ProjectVersion}{}

\newcommand{\Proponente}{Zucchetti S.p.A.}
\newcommand{\Committente}{Prof. Tullio Vardanega \\ Prof. Riccardo Cardin}
\newcommand{\CommittenteInline}{Prof. Tullio Vardanega, Prof. Riccardo Cardin}
\newcommand{\ResponsabileInCarica}{\mattia}

% La versione dei documenti deve essere definita qui in global, perchè serve anche agli altri documenti
\newcommand{\VersioneNP}{1.0.1}  % Norme di Progetto
\newcommand{\VersioneAR}{1.0.2}  % Analisi dei Requisiti
\newcommand{\VersionePQ}{1.0.3}  % Piano di Qualifica
\newcommand{\VersioneG}{1.0.1}   % Glossario
\newcommand{\VersionePP}{1.0.4}  % Piano di Progetto
\newcommand{\VersioneST}{0.0.7} % Specifica Tecnica
\newcommand{\VersioneMU}{0.0.6} % Manuale Utente

% Il verbale non ha versionamento.
% Lasciare vuoto, non mettere trattini o puntini
% Non sono permessi numeri nel nome di un comando :(
\newcommand{\VersioneVprimo}{}
\newcommand{\VersioneVsecondo}{}

% Per riferirsi ad alti documenti globalmente, se la versione nella parte soprastante è vuota il numero non compare
\newcommand{\PianoDiProgetto}{\emph{Piano di Progetto v\VersionePP{}}}
\newcommand{\NormeDiProgetto}{\emph{Norme di Progetto v\VersioneNP{}}}
\newcommand{\PianoDiQualifica}{\emph{Piano di Qualifica v\VersionePQ{}}}
\newcommand{\AnalisiDeiRequisiti}{\emph{Analisi dei Requisiti v\VersioneAR{}}}
\newcommand{\Glossario}{\emph{Glossario v\VersioneG{}}}
\newcommand{\SpecificaTecnica}{\emph{Specifica Tecnica v\VersioneST{}}}
\newcommand{\ManualeUtente}{\emph{Manuale Utente v\VersioneMU{}}}

% Per riferirsi alle revisioni usare i comandi seguenti:
\newcommand{\RTB}{Requirements and Technology Baseline}
\newcommand{\PB}{Product Baseline}
\newcommand{\CA}{Customer Acceptance}

% Comandi utili per riferirsi alle varie fasi
\newcommand{\AReq}{Analisi dei Requisiti}
\newcommand{\ARischi}{Analisi dei rischi}
\newcommand{\AD}{Analisi di Dettaglio}
\newcommand{\PA}{Progettazione Architetturale}
\newcommand{\PD}{Progettazione di Dettaglio}
\newcommand{\Cod}{Codifica}
\newcommand{\VV}{Verifica e Validazione}

%Comandi shortcut per riferirsi a documenti SENZA versione
\newcommand{\PdP}{\emph{Piano di Progetto}}
\newcommand{\PdQ}{\emph{Piano di Qualifica}}
\newcommand{\NdP}{\emph{Norme di Progetto}}
\newcommand{\AdR}{\emph{Analisi dei Requisiti}}
\newcommand{\Gls}{\emph{Glossario}}
\newcommand{\ST}{\emph{Specifica Tecnica}}
\newcommand{\MU}{\emph{Manuale Utente}}

%Comando shortcut per riferirsi al Way of Working
\newcommand{\WoW}{Way of Working}

% Comandi utili per riferirsi ai componenti del gruppo (Diari delle modifiche)
\newcommand{\tommaso}{Tommaso Stocco}
\newcommand{\sebastiano}{Sebastiano Lewental}
\newcommand{\marco}{Marco Cristo}
\newcommand{\riccardo}{Riccardo Cavalli}
\newcommand{\raul}{Raul Pianon}
\newcommand{\martina}{Martina Dall'Amico}
\newcommand{\mattia}{Mattia Zecchinato}


% Comandi utili per riferirsi ai ruoli di progetto
%% non rimuovere i caratteri % gestiscono gli spazi
\NewDocumentCommand{\Responsabile}{o}{%
  \IfNoValueTF{#1}
    {\textit{responsabile di progetto}}
    {\textit{Responsabile di progetto}}%
}
\NewDocumentCommand{\Amministratore}{o}{%
  \IfNoValueTF{#1}
    {\textit{amministratore}}
    {\textit{Amministratore}}%
}
\NewDocumentCommand{\Analista}{o}{%
  \IfNoValueTF{#1}
    {\textit{analista}}
    {\textit{Analista}}%
}
\NewDocumentCommand{\Progettista}{o}{%
  \IfNoValueTF{#1}
    {\textit{progettista}}
    {\textit{Progettista}}%
}
\NewDocumentCommand{\Programmatore}{o}{%
  \IfNoValueTF{#1}
    {\textit{programmatore}}
    {\textit{Programmatore}}%
}
\NewDocumentCommand{\Verificatore}{o}{%
  \IfNoValueTF{#1}
    {\textit{verificatore}}
    {\textit{Verificatore}}%
}
\NewDocumentCommand{\Redattore}{o}{%
  \IfNoValueTF{#1}
    {\textit{redattore}}
    {\textit{Redattore}}%
}

\NewDocumentCommand{\Amministratrice}{o}{%
  \IfNoValueTF{#1}
    {\textit{amministratrice}}
    {\textit{Amministratrice}}%
}
\NewDocumentCommand{\Programmatrice}{o}{%
  \IfNoValueTF{#1}
    {\textit{programmatrice}}
    {\textit{Programmatrice}}%
}
\NewDocumentCommand{\Redattrice}{o}{%
  \IfNoValueTF{#1}
    {\textit{redattrice}}
    {\textit{Redattrice}}%
}

% FIXME: rimuovere
% \newcommand{\Responsabile}{\textit{responsabile di progetto}}
% \newcommand{\Amministratore}{\textit{amministratore}}
% \newcommand{\Analista}{\textit{analista}}
% \newcommand{\Progettista}{\textit{progettista}}
% \newcommand{\Programmatore}{\textit{programmatore}}
% \newcommand{\Verificatore}{\textit{verificatore}}
% \newcommand{\Redattore}{\textit{redattore}}

% \newcommand{\Amministratrice}{\textit{amministratrice}}
% \newcommand{\Programmatrice}{\textit{programmatrice}}
% \newcommand{\Redattrice}{\textit{redattrice}}


\newcommand{\ScopoDelProdotto}{}

\newcommand{\ScopoDelProdottoEng}{}

\newcommand{\GlossarioIntroduzione}{
  Allo scopo di evitare incomprensioni relative al linguaggio utilizzato nella documentazione di progetto, viene fornito un \Gls, nel quale ciascun termine è corredato da una spiegazione che mira a disambiguare il suo significato. I termini tecnici, gli acronimi e i vocaboli ritenuti ambigui vengono formattati in corsivo all'interno dei rispettivi documenti e marcati con una lettera \ped{G} in pedice. Tutte le ricorrenze di un termine definito nel \Gls\ subiscono la formattazione sopracitata.
}

%Comando per i verbali esterni (solamente quando ci sono tanti termini tecnici o ambigui)
\newcommand{\GlossarioIntroduzioneVE}{
  Allo scopo di evitare incomprensioni relative al linguaggio utilizzato nella documentazione di progetto, viene fornito un \Gls, nel quale ciascun termine è corredato da una spiegazione che mira a disambiguare il suo significato. I termini tecnici, gli acronimi e i vocaboli ritenuti ambigui vengono formattati in corsivo all'interno dei rispettivi documenti e marcati con una lettera \ped{G} in pedice. In questo documento viene formattata solamente la prima ricorrenza di un termine definito nel \Gls.
}

\newcommand{\ToBeContinued}{\textit{Continua nella prossima pagina}}

% Altri comandi utili
\newcommand{\email}{email}
\newcommand*{\sezione}[1]{\hyperref[{#1}]{sezione §\ref*{#1}}}
\newcommand*{\appendice}[1]{\hyperref[{#1}]{Appendice \autoref*{#1}}}

% Bisognerebbe aggiungere anche altri brand utilizzati spesso nei vari documenti
\newcommand{\Macos}{macOs}

%%%%%%%%%%%%%%
%  FUNZIONI  %
%%%%%%%%%%%%%%

% In questa seconda parte vanno definite le 'funzioni' utilizzate da due o più documenti.

% Comando per usare comodamente le virgolette
\newcommand{\virgolette}[1]{``#1''}

\newcommand\flextt{%
	\fontdimen2\font=0.2em% interword space
	\fontdimen3\font=0.1em% interword stretch
	\fontdimen4\font=0.1em% interword shrink
	\fontdimen7\font=0.2em% extra space
	\hyphenchar\font=`\-
}

% Serve a dare la giusta formattazione al codice inline
\newcommand{\code}[1]{\ttfamily\flextt{#1}\normalfont}

% Serve a dare la giusta formattazione a tutte le path presenti nei documenti
\newcommand{\file}[1]{\flextt{#1}}

% Permette di andare a capo all'interno di una cella in una tabella
\newcommand{\multiLineCell}[2][c]{\begin{tabular}[#1]{@{}l@{}}#2\end{tabular}}

% Per l'allineamento delle immagini
\usepackage[export]{adjustbox}

% Genera automaticamente la pagina di copertina
\newcommand{\makeFrontPage}{
  % Declare new goemetry for the title page only.
  \newgeometry{top=2cm}
  \begin{titlepage}
  \begin{center}
  \vspace*{2cm}
  \begin{center}
  \includegraphics[width=8cm]{argo_square.png}
  \end{center}

%  \vspace{1cm}

  \begin{Huge}
  \textbf{\DocTitle{}}
  \end{Huge}

  \textbf{\emph{Gruppo \GroupName{} \, \texttwelveudash{} \, Progetto \ProjectName{}}}

  \vspace{10pt}

  \bgroup
  \def\arraystretch{1.3}
  \begin{tabular}{ r|l }
    \multicolumn{2}{c}{\textbf{Informazioni sul documento}} \\
    \hline
    \ifx\DocVersion\empty
      \ifdefined\DocData
        \ifx\DocData\empty
          % lascio vuoto
        \else
          \textbf{Data} & \DocData{} \\
        \fi
      \fi
    \else
      \textbf{Versione} & \faIcon{tag} \DocVersion{} \\
    \fi
    \textbf{Approvazione} & \multiLineCell[t]{\DocApprovazione{}} \\
    \textbf{Uso} & \DocUso{} \\
    \textbf{Distribuzione} & \multiLineCell[t]{\DocDistribuzione{}} \\
  \end{tabular}
  \egroup

  \vspace*{\fill}

  \begin{figure}[H]
	\includegraphics[height=3.5cm, center]{\ModelAssets{logo_unipd.png}}
  \end{figure}
  \end{center}
  \end{titlepage}

  % Ends the declared geometry for the titlepage
  \restoregeometry
}

%%%%%%%%%%%%%%
%  COSTANTI  %
%%%%%%%%%%%%%%

% In questa prima parte vanno definite le 'costanti' utilizzate soltanto da questo documento.
% Devono iniziare con una lettera maiuscola per distinguersi dalle funzioni.

\newcommand{\DocTitle}{Verbale Riunione 2024-05-02}
\newcommand{\DocVersion}{1.0.0}
\newcommand{\DocApprovazione}{\raul}
\newcommand{\DocUso}{Interno}
\newcommand{\DocDistribuzione}{
	\Committente{} \\
	Gruppo \GroupName{}
}

% La descrizione del documento
\newcommand{\DocDescription}{
	Durante l'incontro sono state discusse le attività svolte durante la settimana e sono stati presentati e discussi i punti per il diario di bordo del 3 maggio 2024. È stata inoltre inviata una mail alla Proponente e fissato un incontro per il 6 maggio 2024.
}

%%%%%%%%%%%%%%
%  FUNZIONI  %
%%%%%%%%%%%%%%

% In questa seconda parte vanno definite le 'funzioni' utilizzate soltanto da questo documento.

\begin{document}

\makeFrontPage

\section*{Registro delle modifiche}

\bgroup
\begin{adjustwidth}{-0.5cm}{-0.5cm}
	% MAX 12.5cm
 	\begin{longtable}{|P{1cm}|P{2.2cm}|P{2.7cm}|P{2.7cm}|>{\arraybackslash}P{3.8cm}|}
	  \hline
		\textbf{Ver.} & \textbf{Data} & \textbf{Redazione} & \textbf{Verifica} & \textbf{Descrizione} \\ 
		\hline
		\endfirsthead

		\hline
		\textbf{Ver.} & \textbf{Data} & \textbf{Redazione} & \textbf{Verifica} & \textbf{Descrizione} \\ 
		\hline
		\endhead

		\hline
		\multicolumn{5}{|r|}{{Continua nella prossima pagina}} \\ 
		\hline
		\endfoot

		\hline
		\endlastfoot

		% IN ORDINE DALLA MODIFICA PIÙ RECENTE ALLA PIÙ VECCHIA
		0.0.1 & 2024-06-09 & \martina & \riccardo & Stesura del documento \\
	\end{longtable}
\end{adjustwidth}
\egroup
\clearpage

\tableofcontents
\clearpage

\section{Introduzione}
Il presente documento include un'analisi dei tre Capitolati esterni disponibili, nei quali le aziende ricoprono il ruolo di Proponente, il Professor Tullio Vardanega e il Professor Riccardo Cardin agiscono in qualità di committenti, mentre il gruppo Argo assume il ruolo di fornitore. Per ciascun Capitolato sono fornite le informazioni generali, l’obiettivo, il dominio applicativo e tecnologico, il cambio di visione delle aziende (in caso si sia verificato nell’arco di questi mesi), e infine un’analisi delle valutazioni personali del team. Soppesando gli aspetti positivi e quelli negativi emersi dagli incontri con le aziende, verranno esposte le motivazioni che hanno portato il gruppo a scartare due opzioni a favore una della terza. 

\section{Valutazione del capitolato scelto}

\subsection{Capitolato C9 - ChatSQL: creare frasi SQL da linguaggio naturale}

\subsubsection{Descrizione}
\begin{itemize}
  \item \textbf{Nome:} ChatSQL
  \item \textbf{Proponente:} Zucchetti S.p.A
  \item \textbf{Committenti:} Prof. Tullio Vardanega e Prof. Riccardo Cardin
\end{itemize}

\subsubsection{Obiettivo}
L’obiettivo del capitolato è sviluppare un software che permetta di generare query SQL valide partendo da richieste in linguaggio naturale utilizzando LLM. 

\subsubsection{Dominio applicativo}
Spinta dalla crescente integrazione di LLM nella vita quotidiana, l’azienda propone un capitolato dove l’uso dell’intelligenza artificiale si fonde con il linguaggio SQL. Si vuole quindi quindi creare un metodo per interrogare database usando il linguaggio naturale. Attualmente, sia ChatGPT che altri sistemi simili richiedono la conoscenza delle tabelle, ma per database di grandi dimensioni, questo approccio risulta difficilmente applicabile. La novità del software proposto è quella di permettere a utenti non familiari con il linguaggio SQL di ottenere query adatte in base al database, interfacciandosi tramite il linguaggio naturale. \\
Affinché il sistema funzioni, un attore che conosce il database deve configurare il software caricando il dizionario dati.\\
\\ 
Le richieste minime del capitolato sono le seguenti:
\begin{itemize}
  \item Archiviazione della descrizione della struttura di un database;
  \item Realizzazione di un’interfaccia per l’inserimento di una frase di interrogazione del database in linguaggio naturale;
  \item Creazione di una procedura che combini i due punti precedenti creando un “prompt” ottimale (che contiene solo le parti di database rilevanti), per interrogare un LLM e ottenere la query in SQL che soddisfi la richiesta;
  \item Sviluppo di un sistema di integrazione che coinvolga i tre punti precedenti.
\end{itemize}
I punti cruciali riguardano la formulazione di un dizionario dati appropriato e l’estrazione delle tabelle coinvolte nella richiesta. Il tutto in modo efficiente e integrato come descritto nel punto 3. 

\subsubsection{Dominio tecnologico}
La Proponente lascia libera scelta sulle tecnologie, dando alcuni consigli:
\begin{itemize}
  \item Python, CSS e JavaScript per lo sviluppo;
  \item Libreria txtai per la ricerca semantica ed esempi di fine-tuning; 
  \item Piattaforma di confronto tra LLM (Chatbot Arena LMSYS);
  \item JSON, XML per semplificare la descrizione del contenuto del database;
  \item ChatGPT o altri modelli per generazione delle query;
  \item LLaMA C++ come risposta alle richieste di hardware sofisticato da parte dei LLM;
  \item API di ChatGPT per una connessione diretta ai LLM;
  \item LM Studio per accedere a un server, con le credenziali di OpenAI, direttamente sulla macchina locale;
  \item GitHub come sistema di versionamento e ITS (Issue tracking system);
  \item Documentazione online per fine-tuning di modelli esistenti;
  \item Hugging Face per la ricerca e la selezione dei modelli LLM.
\end{itemize}

\subsection{Mutamento visione acquirente negli ultimi mesi}
Non è stato evidenziato nessun cambiamento significativo in seguito al lavoro svolto dai precedenti fornitori.

\subsubsection{Valutazione}
\begin{itemize}
  \item Aspetti positivi:
  \begin{itemize}
    \item Ampia disponibilità di materiale online (facilmente reperibile) su LLM e SQL;
    \item Libertà sul dominio tecnologico;
    \item L’idea ha suscitato interesse nel gruppo poiché si tratta di lavorare con tecnologie moderne e innovative;
    \item L’obiettivo del capitolato rappresenta una sfida stimolante e prospetta l’apprendimento di  utili in futuro;
    \item A seguito del lavoro dei precedenti fornitori, la Proponente ha espresso maggior sicurezza sulla fattibilità dei requisiti opzionali;
    \item Nessun componente del gruppo ha esperienze pregresse sugli argomenti del capitolato e pertanto questo si presenta come una sfida stimolante.
  \end{itemize}
  \item Aspetti negativi:
  \begin{itemize}
    \item L’approccio alla costruzione del dizionario dati è poco chiaro;
    \item La parte di pulizia pre-prompt potrebbe risultare complessa;
    \item Ripetitività del processo di prompt engineering, vista la sua natura iterativa basata su trial-and-error.
  \end{itemize}
\end{itemize}

\subsubsection{Conclusioni:}
Il capitolato ha suscitato interesse nel gruppo poichè prevede l'utilizzo di strumenti moderni e all'avanguardia, lasciando però al team la libertà sul dominio tecnologico. Nonostante alcuni fattori potrebbero rappresentare degli ostacoli non indifferenti, l'esperienza e le competenze acquisite sono considerate fondamentali per il percorso accademico e il futuro lavorativo del team. Inoltre, la disponibilità dell'azienda a fornire materiali e sostegno ha rafforzato la posizione del gruppo sul capitolato.

\section{Valutazione capitolati rimanenti}

\subsection{Capitolato C6 - SyncCity: Smart city monitoring platform}

\subsubsection{Descrizione}
\begin{itemize}
  \item \textbf{Nome:} SyncCity
  \item \textbf{Proponente:} SyncLab S.r.l.
  \item \textbf{Committenti:} Prof. Tullio Vardanega e Prof. Riccardo Cardin
\end{itemize}

\subsubsection{Obiettivo}
L'obiettivo del capitolato è sviluppare una piattaforma di monitoraggio di smart city che sia in grado di acquisire dati da dispositivi IoT o sensori, per poi elaborarli e renderli disponibili agli utenti mediante una dashboard.

\subsubsection{Dominio applicativo}
La Proponente si focalizza sullo sviluppo, con particolare attenzione al lato back-end, di una piattaforma di monitoraggio ambientale applicata al contesto di una smart city, con la successiva possibile verticalizzazione nell'ambito dello smart parking. L'enfasi viene dunque posta sulla data integration, che dovrà essere affiancata da un'appropriata simulazione di dati provenienti da numerose fonti, come ad esempio:
\begin{itemize}
  \item Sensori di temperatura;
  \item Sensori di polveri sottili;
  \item Sensori di umidità;
  \item Sensori di guasti elettrici.
\end{itemize}

Di conseguenza, la soluzione richiesta prevede lo sviluppo di un sistema di generazione dei dati mediante la simulazione di sensori, l’inserimento delle informazioni in una coda e la loro successiva elaborazione. L'utente finale potrà consultare i dati tramite una piattaforma di data visualization. Quest'ultima rappresenta la componente front-end dell'applicativo, e prevede la creazione di una maschera contenente la mappa della città con le informazioni raccolte.

\subsubsection{Dominio tecnologico}
La Proponente suggerisce le seguenti tecnologie, considerate strumenti adatti per operazioni come il data stream processing e la gestione di Big Data:
\begin{itemize}
  \item Script Python ed eventualmente librerie di generazione dati come faker per lo sviluppo di un framework per la simulazione dei dati;
  \item Apache Kafka come broker per disaccoppiare lo stream di informazioni provenienti dai simulatori;
  \item ClickHouse come database OLAP colonnare, per la sua capacità di gestire grandi moli di dati ad alta velocità e per la sua integrazione agevole con Kafka;
  \item Docker per la containerizzazione;
  \item Grafana come piattaforma di data visualization.
\end{itemize}
L’azienda resta comunque disponibile a considerare eventuali variazioni richieste dal fornitore, purché soddisfino le necessità del capitolato.

\subsection{Mutamento visione acquirente negli ultimi mesi}
A seguito della ricezione di prodotti sufficientemente aderenti alle richieste, la Proponente sta valutando la possibilità (se applicabile) di ampliare i requisiti oltre il Proof Of Concept, per approfondire l'ambito dello smart parking.\\
Di seguito ne sono elencati gli aspetti chiave:
\begin{itemize}
  \item Gestione di entrate e uscite dal parcheggio;
  \item Acquisizione di parametri ambientali;
  \item Transazioni per sosta;
  \item Controllo real-time della correttezza del pagamento sulla base del tempo di 	parcheggio.
\end{itemize}
Lo smart parking costituisce una verticalizzazione della gestione di una smart city e ne ingloba le componenti più concrete.

\subsubsection{Valutazione}
\begin{itemize}
  \item Aspetti positivi:
  \begin{itemize}
    \item Disponibilità della Proponente a organizzare incontri frequenti per SAL (Stato Avanzamento Lavori) o per altre necessità di progetto, anche in sede;
    \item Disponibilità della Proponente a fornire supporto asincrono tramite l’accesso a una chat per la comunicazione diretta con i dipendenti dell'azienda;
    \item Disponibilità per deep dives sulle tecnologie consigliate per lo svolgimento del 	capitolato;
    \item Chiarezza comunicativa della Proponente;
    \item Concretezza degli obiettivi e delle possibilità didattiche;
    \item Interesse del gruppo sulle tematiche di IoT e Big Data;
    \item Il gruppo sarebbe proprietario dell'applicativo al termine del progetto;
    \item Conoscenza pregressa dell'ambito di applicazione e di alcune delle tecnologie proposte.
  \end{itemize}
  \item Aspetti negativi:
  \begin{itemize}
    \item L'utilizzo di nuove tecnologie, differenti da quelle approfondite durante il percorso di studi, comporta una notevole curva di apprendimento;
    \item Gli obiettivi successivi alla consegna minima si discostano parecchio da quest'ultima in termini di complessità e tempistiche;
    \item La gestione di dati con formati differenti e la configurazione Kafka - ClickHouse potrebbero risultare operazioni dispendiose.
  \end{itemize}
\end{itemize}

\subsubsection{Conclusioni}
Questo progetto risulta avere molti aspetti positivi tuttavia, una conoscenza superficiale delle tecnologie e dell’argomento ha portato il gruppo ad allontanarsi da quest’ultimo.\\
L’interesse era comunque forte, ma vista la dimensione del progetto e la quantità di tempo che avrebbe richiesto, al termine di diversi incontri, questo è stato selezionato come seconda scelta.

\subsection{Capitolato C3 - EasyMeal: Prenota, Ordina, Mangia}

\subsubsection{Descrizione}
\begin{itemize}
  \item \textbf{Nome:} EasyMeal
  \item \textbf{Proponente:} Imola informatica
  \item \textbf{Committenti:} Prof. Tullio Vardanega e Prof. Riccardo Cardin
\end{itemize}

\subsubsection{Obiettivo}
L'obiettivo del capitolato è sviluppare una web app responsive che sia in grado di affrontare alcune problematiche legate al settore dei ristoranti; in particolare, la complessità delle prenotazioni e l'inefficienza nelle ordinazioni. EasyMeal mira a rendere la prenotazione di un tavolo un'esperienza coinvolgente e collaborativa, instaurando un canale di comunicazione tra gli utenti e lo staff. Inoltre, i clienti possono pianificare l’orario di arrivo al ristorante e creare il proprio ordine da casa o dall'ufficio. Oltre a semplificare il processo di prenotazione e ordinazione dal punto di vista dei clienti, l’applicazione deve agevolare la pianificazione della spesa a carico dei ristoratori, per ridurre lo spreco di risorse alimentari.

\subsubsection{Dominio applicativo}
EasyMeal punta a migliorare l’efficienza, la personalizzazione e la comodità delle prenotazioni e delle ordinazioni nel settore dei ristoranti, offrendo un servizio di qualità superiore sia ai clienti che ai ristoratori.
\\ 
L’applicazione mette a disposizione le seguenti funzionalità:
\begin{itemize}
  \item Registrazione di un nuovo utente (base o amministratore);
  \item Sistema di filtraggio dei ristoranti per data, orario, città e tipologia di cucina;
  \item Indicazione di eventuali allergie e intolleranze in fase di registrazione;
  \item Prenotazione di un tavolo;
  \item Sistema di richiesta/conferma della prenotazione tra cliente e ristoratore;
  \item Ordinazione collaborativa dei pasti;
  \item Interazione con lo staff del ristorante tramite una chat privata;
  \item Divisione (equa o proporzionale) del conto;
  \item Consultazione real-time delle prenotazioni da parte di un amministratore del ristorante;
  \item Monitoraggio continuo dei pagamenti da parte di un amministratore del ristorante;
  \item Rilascio di feedback e recensioni sul servizio complessivo.
\end{itemize}

\subsubsection{Dominio tecnologico}
Trattandosi di un’applicazione riservata non solo ai clienti, ma anche allo staff dei ristoranti, EasyMeal deve essere accessibile sia da PC che da mobile (IOS e Android). Non ci sono vincoli sulle tecnologie da utilizzare, tuttavia gli strumenti comunemente impiegati per sviluppare web app sono i seguenti:
\begin{itemize}
  \item Front-end: linguaggi come JavaScript, HTML, CSS;
  \item Front-end: framework e librerie come React.js, Angular, Vue.js;
  \item Back-end: linguaggi come Python, Java, PHP;
  \item Back-end: framework come Express.js, Flask, Laravel;
  \item Database: MySQL, PostgreSQL, MongoDB, SQLite.
\end{itemize}

\subsubsection{Valutazione}
\begin{itemize}
  \item Aspetti positivi:
  \begin{itemize}
    \item Buona definizione e presentazione dei requisiti minimi;
    \item Libertà sul dominio tecnologico;
    \item Inserimento (nel documento di presentazione) di un glossario con alcune definizioni per approcciarsi al progetto EasyMeal;
    \item L'applicazione si colloca precisamente in un contesto reale e moderno;
    \item Il capitolato ha suscitato interesse per la concretezza del progetto, che parte da un’accezione semplice come quella dell’app per i ristoranti, ma integra al suo interno degli aspetti originali e innovativi;
    \item La relativa semplicità delle tecnologie e delle funzionalità richieste, già affrontate in passato da alcuni componenti del gruppo;
    \item L'azienda mette a disposizione macchine virtuali nelle quali effettuare le installazioni dei componenti applicativi;
    \item L’azienda mette a disposizione professionisti per fornire supporto dal punto di vista tecnico e architetturale.
  \end{itemize}
  \item Aspetti negativi:
  \begin{itemize}
    \item La quantità di requisiti (minimi e opzionali) preoccupa per l’effettiva lunghezza del progetto;
    \item Parecchi requisiti opzionali sono visti come forzati (in particolare la richiesta di cifrare i dati sensibili);
    \item La disponibilità dell’azienda sembra essere meno costante rispetto alle Proponenti di altri capitolati;
    \item L’ambiguità nelle tecnologie e nell’architettura;
    \item La struttura base e minimale del prodotto è risultata poco chiara;
    \item Il capitolato richiede l'esecuzione di tante operazioni relativamente semplici e meccaniche (per esempio inserimenti, modifiche, aggiornamenti, rimozioni), a cui vanno integrate funzionalità decisamente più complesse ed elusive;
    \item Diversi componenti del gruppo hanno già lavorato con le tecnologie e gli strumenti richiesti dal capitolato, che non rappresentano quindi una sfida di apprendimento stimolante.
  \end{itemize}
\end{itemize}

\subsubsection{Conclusioni}
Dopo una scrupolosa valutazione del capitolato, soppesando punti di forza e criticità, il gruppo ha deciso di classificare come terza scelta questo progetto. In particolare, il team ha evidenziato una buona definizione dei requisiti minimi, i quali però non sono risultati abbastanza accattivanti. Inoltre, l’abbondanza delle richieste (considerando anche la quasi obbligatorietà di alcuni requisiti opzionali) rischia di tradursi in un tempo di sviluppo eccessivo. L’idea, seppur valida e concreta, è stata considerata meno affascinante rispetto agli altri due capitolati presi in considerazione, i quali spiccano per innovazione, contemporaneità e perciò richiedono l’impiego di tecnologie all’avanguardia.

\begin{samepage}

  \vspace*{\fill}
  Luogo e Data: \\
  Padova (PD) 2024-06-22
  \vspace*{18pt}

  \begin{tikzpicture}[overlay]
    \node [anchor = south west] at (0cm,-0.14cm) {Firma:};
    % Firma interna da sistemare
    %\node [anchor = south east] at (\textwidth- -1cm,-1.3cm) {\includegraphics[width=8cm]{\ModelAssets{signatures/firma_sebastiano.png}}}; 
    \draw [anchor = south west] (1.5cm,0cm) -- (\columnwidth,0cm);
  \end{tikzpicture}
  
  \vspace*{20pt}
\end{samepage}



\end{document}
